% chktex-file 36

\documentclass[12pt]{report}

% chktex-file 1

\usepackage{graphicx} % Required for inserting images
\usepackage[utf8]{inputenc}
\usepackage[a4paper]{geometry} % Adjust margins
\usepackage{setspace} % Used for adding spaces in contents
\usepackage{amsfonts,amsmath,amssymb,amsthm}
\usepackage{enumitem}
\usepackage{footnote}
\usepackage{siunitx} % Writing SI units
\usepackage[titles]{tocloft} % Table of contents titles made bold
\usepackage[labelfont=bf]{caption} % Boldens names of figures
\usepackage[bottom]{footmisc} % Footnotes at bottom of page
\usepackage[hidelinks]{hyperref}
\usepackage[all]{hypcap} % Hyperlink sends you to top of reference

\usepackage[
    backend=biber,
    citestyle=authoryear,
    sorting=ynt,
]{biblatex} % BibTeX for bibliography

\hypersetup{
    colorlinks=true,
    linkcolor=blue,,
    urlcolor=blue,
    citecolor=blue
}

\renewcommand{\labelitemi}{\(\bullet\)}
\renewcommand{\labelitemii}{\(\circ\)}

\newcommand\partialdev[2]{\frac{\partial #1}{\partial #2}}
\newcommand\dev[2]{\frac{d #1}{d #2}}
\newcommand\codev[2]{\frac{D #1}{D #2}}

\setlist[itemize]{itemsep=1pt,topsep=3pt}
\setlist[enumerate]{label=(\roman*{}),itemsep=1pt,topsep=3pt}
\newcommand{\boldu}{\mathbf{u}}

\title{\textbf{Oceanic Mixed Bottom Layers}}
\author{Xing Yu Yi}
\date{August 2026}

\addbibresource{bibliography.bib}

\begin{document}

\maketitle

\hypersetup{
    linkcolor=black,
}
\tableofcontents
\hypersetup{
    linkcolor=blue,
}

\chapter{Introduction}
\input{Chapters/Introduction}\label{intro}

\chapter{Ekman model}\label{chap:2}

We will begin by looking at an Ekman model for the formation of the bottom boundary layer, where we will have the effects of the Coriolis force with a far-stream velocity \(\mathbf{U}_\infty\). The flow will initially be in geostrophic balance, with a pressure gradient to balance this. This model will largely be based off~\cite{TaylorSarkar2008} and looking at stratified Ekman layers.

\input{Chapters/Ekman}

\chapter{Stokes model}\label{chap:3}

We will next look at the Stokes model

\printbibliography{}

\end{document}